\documentclass[11pt]{article}
\usepackage{mathblatt}
\makeatletter
\setlength{\mbpfbe}{0mm}% keine Punkte (Bauregel 6.4)
\makeatother
% Rechenraum als Karo (Bauregel 6.9): n Rechenschritte = n mal 10 mm
\newcommand{\karo}[1]{\par\smallskip\noindent\begin{tikzpicture}
  \clip (0,0) rectangle ({\linewidth-1mm},{#1*10mm});
  \draw[black!16,line width=0.3pt,step=5mm] (0,0) grid ({\linewidth},{#1*10mm});
  \end{tikzpicture}\par}
\newcommand{\linie}[1]{\,{\color{black!45}\rule[-1pt]{#1}{0.4pt}}\,}
\newcommand{\pkt}{\textperiodcentered{}}
\begin{document}

\pfheftstil
\pflaufkopf[P10 \pkt{} mehr]{Baumdiagramm und Pfadregeln \pkt{} Lösungen}{PS50-2}\pfseitevon
\pfheftkopf{Baumdiagramm und Pfadregeln}{Lösungen \pkt{} mehr}
\begin{pfloesungb}{}
\lzb{1.}{}{P $\approx$ 16,7 \%}{(1,4), (2,4), (3,4), (4,4), (5,4), (6,4); P = $\tfrac{6}{36}$ = $\tfrac{1}{6}$ $\approx$ 16,7 \%}
\end{pfloesungb}
\begin{pfloesungb}{}
\lzb{2.}{}{$\tfrac{4}{24}$ $\approx$ 16,7 \%}{$\tfrac{4}{24}$ = $\tfrac{1}{6}$ $\approx$ 16,7 \%}
\end{pfloesungb}
\begin{pfloesungb}{}
\lzb{3.}{}{8 Ergebnisse (ZZZ, ZZW, ZWZ, ZWW, WZZ, WZW, WWZ, WWW)}{8 Ergebnisse (ZZZ, ZZW, ZWZ, ZWW, WZZ, WZW, WWZ, WWW)}
\end{pfloesungb}
\begin{pfloesungb}{}
\lzb{4.}{}{$\tfrac{1}{8}$ = 12,5 \%}{$\tfrac{1}{8}$ = 12,5 \%}
\end{pfloesungb}
\begin{pfloesungb}{}
\lzb{5.}{}{$\tfrac{3}{50}$ = 6 \%}{$\tfrac{3}{50}$ = 6 \%}
\end{pfloesungb}
\begin{pfloesungb}{}
\lzb{6.}{}{$\tfrac{1}{600}$}{700 $-$ 101 + 1 = 600 Lose, 1 Hauptgewinn $\Rightarrow$ $\tfrac{1}{600}$}
\end{pfloesungb}
\begin{pfloesungb}{}
\lzb{7.}{}{3 rote Felder}{3 rote Felder}
\end{pfloesungb}
\begin{pfloesungb}{}
\lzb{8.}{}{2. Stufe an jedem Ast: A 4/10, B 3/10, C $\tfrac{3}{10}$}{1. Stufe: A 4/10, B $\tfrac{3}{10}$; 2. Stufe an jedem Ast: A 4/10, B 3/10, C $\tfrac{3}{10}$}
\end{pfloesungb}
\begin{pfloesungb}{}
\lzb{9.}{}{mittlerer Topf (2 weiße, 6 graue Kugeln): $\tfrac{2}{8}$ = 25 \%}{mittlerer Topf (2 weiße, 6 graue Kugeln): $\tfrac{2}{8}$ = 25 \%}
\end{pfloesungb}
\begin{pfloesungb}{}
\lzb{10.}{}{P(Gewinn) = 25 \%}{P(Gewinn) = $\tfrac{30}{120}$ = $\tfrac{1}{4}$ = 25 \%}
\end{pfloesungb}
\begin{pfloesungb}{}
\lzb{11.}{}{Max hat nicht Recht, beide Wahrscheinlichkeiten sind $\tfrac{1}{6}$ $\approx$ 0,167}{Max hat nicht Recht, beide Wahrscheinlichkeiten sind $\tfrac{1}{6}$ $\approx$ 0,167}
\end{pfloesungb}
\begin{pfloesungb}{}
\lzb{12.}{}{$\tfrac{4}{6}$ $\approx$ 0,67}{$\tfrac{4}{6}$ = $\tfrac{2}{3}$ $\approx$ 0,67}
\end{pfloesungb}
\begin{pfloesungb}{}
\lzb{13.}{}{4 weiße Kugeln}{4 weiße Kugeln}
\end{pfloesungb}
\begin{pfloesungb}{}
\lzb{14.}{}{$\tfrac{3}{10}$; 30 \%}{$\tfrac{3}{10}$; 30 \%}
\end{pfloesungb}
\begin{pfloesungb}{}
\lzb{15.}{}{nicht rot: $\tfrac{19}{24}$; nicht rot > rot > rot: $\tfrac{4}{22}$; Aussage 1 falsch, Aussage 2 falsch, Aussage 3 wahr (P $\approx$ 0,49 \%)}{nicht rot: $\tfrac{19}{24}$; nicht rot > rot > rot: $\tfrac{4}{22}$; Aussage 1 falsch, Aussage 2 falsch, Aussage 3 wahr (P = $\tfrac{60}{12144}$ $\approx$ 0,49 \%)}
\end{pfloesungb}
\begin{pfloesungb}{}
\lzb{16.}{}{nR 12/13, nR>R 1/12, nR>nR>R 1/11, nR>nR>nR $\tfrac{10}{11}$; P $\approx$ 0,23}{nR 12/13, nR>R 1/12, nR>nR>R 1/11, nR>nR>nR $\tfrac{10}{11}$; P = $\tfrac{3}{13}$ $\approx$ 0,23}
\end{pfloesungb}
\begin{pfloesungb}{}
\lzb{17.}{}{nein, nicht $\tfrac{3}{40}$ $\approx$ 7,9 \%}{nein, nicht $\tfrac{3}{40}$ = 7,5 \%, sondern P = $\tfrac{3}{38}$ $\approx$ 7,9 \%}
\end{pfloesungb}
\begin{pfloesungb}{}
\lzb{18.}{}{$\tfrac{34}{100}$ = 0,34}{$\tfrac{34}{100}$ = $\tfrac{17}{50}$ = 0,34}
\end{pfloesungb}
\begin{pfloesungb}{}
\lzb{19.}{}{$\tfrac{5}{60}$}{ja, 5 Trostpreis-Lose (144, 344, 444, 544, 644) unter 60 Losen mit Endziffer 4 $\Rightarrow$ $\tfrac{5}{60}$}
\end{pfloesungb}
\begin{pfloesungb}{}
\lzb{20.}{}{P $\approx$ 7,4 \%}{alle Äste „4“ 1/6, „keine 4“ $\tfrac{5}{6}$; P = $\tfrac{16}{216}$ = $\tfrac{2}{27}$ $\approx$ 7,4 \%}
\end{pfloesungb}
\begin{pfloesungb}{}
\lzb{21.}{}{P(3 blau) $\approx$ 0,99 \%; P(3 gelb) $\approx$ 0,20 \%; Behauptung falsch, es ist das 5-fache und nicht das Doppelte}{P(3 blau) = $\tfrac{120}{12144}$ = $\tfrac{5}{506}$ $\approx$ 0,99 \%; P(3 gelb) = $\tfrac{24}{12144}$ = $\tfrac{1}{506}$ $\approx$ 0,20 \%; Behauptung falsch, es ist das 5-fache und nicht das Doppelte}
\end{pfloesungb}
\begin{pfloesungb}{}
\lzb{22.}{}{S): 3 Ergebnisse; P $\approx$ 0,854; E: beide Pfannkuchen haben dieselbe Füllung, P(E) $\approx$ 0,858}{(S;M), (M;S), (S;S): 3 Ergebnisse; P = $\tfrac{37}{40}$ · $\tfrac{36}{39}$ = $\tfrac{1332}{1560}$ $\approx$ 0,854; E: beide Pfannkuchen haben dieselbe Füllung, P(E) $\approx$ 0,858}
\end{pfloesungb}
\end{document}